\documentclass[11pt,a4paper]{article}
\usepackage[utf8]{inputenc}
\usepackage[margin=2.2cm,top=2.5cm,bottom=2.5cm]{geometry}
\usepackage{xcolor}
\usepackage{listings}
\usepackage{tcolorbox}
\tcbuselibrary{skins,breakable}
\usepackage{fancyhdr}
\usepackage{hyperref}
\usepackage{booktabs}
\usepackage{tabularx}
\usepackage{enumitem}
\usepackage{microtype}
\usepackage{titlesec}

% Colors
\definecolor{primaryblue}{RGB}{29, 78, 216}
\definecolor{darkslate}{RGB}{15, 23, 42}
\definecolor{terminalbg}{RGB}{15, 23, 42}
\definecolor{terminalfg}{RGB}{248, 250, 252}
\definecolor{terminalborder}{RGB}{30, 41, 59}
\definecolor{conceptborder}{RGB}{37, 99, 235}
\definecolor{conceptbg}{RGB}{239, 246, 255}
\definecolor{noteborder}{RGB}{22, 163, 74}
\definecolor{notebg}{RGB}{240, 253, 244}
\definecolor{lightgray}{RGB}{241, 245, 249}

% Header & Footer
\pagestyle{fancy}
\fancyhf{}
\renewcommand{\headrulewidth}{0.4pt}
\renewcommand{\headrule}{\hbox to\headwidth{\color{lightgray}\leaders\hrule height \headrulewidth\hfill}}
\fancyhead[L]{\small\color{gray}\textbf{RecordiaILIA} \textbar\ Collaborative Git \& GitHub Workflow Report}
\fancyhead[R]{\small\color{gray}November 2026}
\fancyfoot[L]{\small\color{gray}Academic Practical Work (TP) Report}
\fancyfoot[R]{\small\color{gray}Page \thepage}

% Hyperref setup
\hypersetup{
    colorlinks=true,
    linkcolor=primaryblue,
    filecolor=primaryblue,      
    urlcolor=primaryblue,
    pdftitle={RecordiaILIA - Collaborative Git & GitHub Workflow Report},
    pdfauthor={Yassine Ouali, Hachem Squalli}
}

% Section formatting
\titleformat{\section}{\Large\bfseries\color{primaryblue}}{\thesection}{1em}{}
\titleformat{\subsection}{\large\bfseries\color{darkslate}}{\thesubsection}{1em}{}

% Terminal / Code listings style
\lstset{
    backgroundcolor=\color{terminalbg},
    basicstyle=\color{terminalfg}\ttfamily\footnotesize,
    breaklines=true,
    frame=none,
    columns=fullflexible,
    keepspaces=true,
    showstringspaces=false
}

% Custom Tutorial Boxes
\newtcolorbox{terminalbox}[1][]{
    enhanced,
    colback=terminalbg,
    colframe=terminalborder,
    coltext=terminalfg,
    arc=4mm,
    boxrule=1pt,
    left=8pt,
    right=8pt,
    top=8pt,
    bottom=8pt,
    fontupper=\ttfamily\footnotesize,
    breakable,
    #1
}

\newtcolorbox{gitconcept}[1]{
    enhanced,
    colback=conceptbg,
    colframe=conceptborder,
    boxrule=0pt,
    leftrule=4pt,
    arc=2mm,
    title={\textbf{\color{primaryblue}Git Principle: #1}},
    coltitle=primaryblue,
    attach title to upper={\par\vspace{3pt}},
    fontupper=\small,
    top=6pt,
    bottom=6pt,
    left=8pt,
    right=8pt
}

\newtcolorbox{notebox}[1]{
    enhanced,
    colback=notebg,
    colframe=noteborder,
    boxrule=0pt,
    leftrule=4pt,
    arc=2mm,
    title={\textbf{\color{noteborder}#1}},
    coltitle=noteborder,
    attach title to upper={\par\vspace{3pt}},
    fontupper=\small,
    top=6pt,
    bottom=6pt,
    left=8pt,
    right=8pt
}

\begin{document}

% -------------------------------------------------------------
% TITLE & HEADER
% -------------------------------------------------------------
\begin{center}
    {\Huge\bfseries\color{primaryblue} RecordiaILIA}\\[0.4em]
    {\Large\bfseries\color{darkslate} Collaborative Git \& GitHub Workflow Report}\\[0.3em]
    {\color{gray}\textit{Academic Practical Work (TP) -- Feature Branching, Issue Tracking, and Multi-Contributor Integration}}\\[1em]
    \rule{\linewidth}{1.2pt}\\[1em]
\end{center}

\begin{minipage}[t]{0.48\textwidth}
    \textbf{Authors \& Roles:}\\[0.3em]
    \textbf{Yassine Ouali} (\texttt{Ysn-Ir} / \texttt{YSN})\\
    {\footnotesize\color{gray} Architecture, Hotkeys, Storage, Video, Steps Recorder, Test Suite}\\[0.3em]
    \textbf{Hachem Squalli} (\texttt{hachemite})\\
    {\footnotesize\color{gray} Architecture, UI Prototype, Core Capture Logic}
\end{minipage}
\hfill
\begin{minipage}[t]{0.48\textwidth}
    \textbf{Repository Details:}\\[0.3em]
    \textbf{GitHub}: \url{https://github.com/hachemite/recordia-ilia}\\
    \textbf{Branch Model}: Feature Branch Workflow (\texttt{--no-ff})\\
    \textbf{Issues Status}: 100\% Closed (8/8 Issues)\\
    \textbf{Test Status}: 20/20 Tests Passing (100\% Pass Rate)
\end{minipage}

\vspace{1.5em}

% -------------------------------------------------------------
% EXECUTIVE SUMMARY
% -------------------------------------------------------------
\section{Executive Summary}

This academic report provides a rigorous, step-by-step documentation of the collaborative Git and GitHub workflows utilized during the development of \textbf{RecordiaILIA}. Designed as a high-performance desktop utility for automated \textbf{Problem Steps Recording (PSR)}, desktop screen video recording, and custom hotkeys, this project demonstrates real-world software engineering practices.

All \textbf{8 original GitHub issues} were completed, peer-reviewed, and merged using isolated feature branches with non-fast-forward (\texttt{--no-ff}) merges. Furthermore, \textbf{3 major advanced enhancements} were designed and integrated: automated step recording with active window and process detection, event-driven triggers (mouse click and periodic timer), and an interactive step review dialog.

\vspace{0.8em}

% -------------------------------------------------------------
% COLLABORATIVE GIT METHODOLOGY
% -------------------------------------------------------------
\section{Collaborative Git Methodology}

The project adheres strictly to the professional \textbf{Feature Branch Workflow}:

\begin{enumerate}[leftmargin=*]
    \item \textbf{Clean Production Branch (\texttt{main})}: The \texttt{main} branch represents production-ready, verified code. Direct commits without feature branches are restricted.
    \item \textbf{Isolated Feature Branches (\texttt{feature/<name>})}: Every distinct issue and functional enhancement is developed on its own dedicated branch.
    \item \textbf{Traceable Non-Fast-Forward Merges (\texttt{--no-ff})}: Merges into \texttt{main} preserve the distinct branch topology in the Git commit graph rather than flattening commits linearly.
    \item \textbf{Automated GitHub Issue Resolution}: Commit messages and merge headers use standard keywords (\texttt{Fixes \#X}, \texttt{Closes \#X}), triggering GitHub's automated issue tracking to close corresponding issues automatically upon merge.
    \item \textbf{Continuous Verification}: Prior to merging, the full automated test suite is run locally, guaranteeing zero regressions.
\end{enumerate}

\vspace{0.8em}

% -------------------------------------------------------------
% GLOBAL COMMIT GRAPH
% -------------------------------------------------------------
\section{Global Commit Graph \& Topology}

The complete topological tree extracted via \texttt{git log --graph --oneline --decorate --all} demonstrates the branching, parallel commits, and merge bubble structure:

\begin{terminalbox}
*   d94d961 (HEAD -> main, origin/main) Merge branch 'feature/hotkey-config-and-step-editor'
|\  
| * 39f309e feat(ui): add custom hotkey configuration, step description editor
|/  
*   2ed5baa Merge branch 'feature/auto-step-triggers' into main
|\  
| * f106f05 feat(triggers): implement automated step capture triggers for mouse click and timer
|/  
*   bbece0a Merge branch 'feature/steps-recorder' into main
|\  
| * aae535f feat(steps): implement automated steps recorder with window inspection
|/  
*   9f431ab Merge branch 'feature/documentation-and-report' (Closes #8)
|\  
| * 60a428d docs: complete user documentation and Git workflow report
|/  
*   798cf8a Merge branch 'feature/test-features' (Closes #7)
|\  
| * 93805f8 test(core): implement comprehensive unit and integration test suite
|/  
*   1b0f505 Merge branch 'feature/screen-recording' (Closes #5)
|\  
| * ee15286 feat(recorder): implement threaded screen recording engine with OpenCV
|/  
*   d2ab6c1 Merge branch 'feature/screenshot-capture' (Closes #4)
|\  
| * b3864cf feat(capture): implement screenshot capture engine with mss
|/  
*   e49e456 Merge branch 'feature/file-saving-and-settings' (Closes #3, Closes #6)
|\  
| * e923205 feat(storage): implement file saving manager and settings dialog
|/  
*   d0735a3 Merge branch 'feature/architecture-and-hotkeys' (Closes #1)
|\  
| * ecb5f0f docs(git): record branch creation, commit, and push in workflow log
| * 4d4966e feat(core): design project architecture and global hotkeys listener (fixes #1)
|/  
*   f0d65b2 Merge pull request #9 from hachemite/feature/Creating_user_interface
|\  
| * cd887b1 feat(ui): create main window with capture buttons (#1)
| * 7d56871 first commit
|/  
* 81f7e7c Initial commit
\end{terminalbox}

\vspace{0.8em}

% -------------------------------------------------------------
% ISSUES & DELIVERABLES MATRIX
% -------------------------------------------------------------
\section{GitHub Issues \& Branch Matrix}

\begin{table}[h!]
\centering
\small
\begin{tabularx}{\textwidth}{c l c l c c}
\toprule
\textbf{Issue} & \textbf{Feature Description} & \textbf{Assignee} & \textbf{Feature Branch} & \textbf{Commit} & \textbf{Status} \\
\midrule
\#1 & Architecture \& Global Hotkeys & \texttt{Ysn-Ir} & \texttt{feature/architecture-and-hotkeys} & \texttt{4d4966e} & \textbf{\color{noteborder}CLOSED} \\
\#2 & User Interface Creation & \texttt{hachemite} & \texttt{feature/Creating\_user\_interface} & \texttt{cd887b1} & \textbf{\color{noteborder}CLOSED} \\
\#3 & File Saving \& Settings Dialog & \texttt{Ysn-Ir} & \texttt{feature/file-saving-and-settings} & \texttt{e923205} & \textbf{\color{noteborder}CLOSED} \\
\#4 & Screenshot Capture Engine & \texttt{Ysn-Ir} & \texttt{feature/screenshot-capture} & \texttt{b3864cf} & \textbf{\color{noteborder}CLOSED} \\
\#5 & Screen Video Recording Engine & \texttt{Ysn-Ir} & \texttt{feature/screen-recording} & \texttt{ee15286} & \textbf{\color{noteborder}CLOSED} \\
\#6 & File Saving \& Settings Dialog & \textit{(unassigned)} & \texttt{feature/file-saving-and-settings} & \texttt{e923205} & \textbf{\color{noteborder}CLOSED} \\
\#7 & Automated Test Suite & \texttt{Ysn-Ir} & \texttt{feature/test-features} & \texttt{93805f8} & \textbf{\color{noteborder}CLOSED} \\
\#8 & Documentation \& Setup Guide & \textit{(unassigned)} & \texttt{feature/documentation-and-report} & \texttt{60a428d} & \textbf{\color{noteborder}CLOSED} \\
Bonus 1 & Automated Steps Recorder & \texttt{Ysn-Ir} & \texttt{feature/steps-recorder} & \texttt{aae535f} & \textbf{\color{primaryblue}MERGED} \\
Bonus 2 & Click \& Timer Auto Triggers & \texttt{Ysn-Ir} & \texttt{feature/auto-step-triggers} & \texttt{f106f05} & \textbf{\color{primaryblue}MERGED} \\
Bonus 3 & Hotkey Re-binding \& Editor & \texttt{Ysn-Ir} & \texttt{feature/hotkey-config-and-step-editor} & \texttt{39f309e} & \textbf{\color{primaryblue}MERGED} \\
\bottomrule
\end{tabularx}
\caption{Complete inventory of completed and verified project features.}
\end{table}

\vspace{0.8em}

% -------------------------------------------------------------
% DETAILED COMMAND LOG & OUTPUTS
% -------------------------------------------------------------
\section{Detailed Command Log \& Terminal Responses}

\subsection{Step 0: Remote Upstream Synchronization}
Before creating feature branches, local \texttt{main} was synchronized with the remote upstream repository:

\begin{terminalbox}
$ git pull origin main
From https://github.com/hachemite/recordia-ilia
 * branch            main       -> FETCH_HEAD
Already up to date.
\end{terminalbox}

\begin{gitconcept}{Upstream Synchronization}
Pulling before branching guarantees that any new branch diverges from the most recent known stable commit on \texttt{origin/main}, preventing silent merge regressions.
\end{gitconcept}

\subsection{Step 1: Issue \#1 -- Project Architecture \& Global Hotkeys Listener}
\begin{itemize}[noitemsep]
    \item \textbf{Goal}: Establish modular Python architecture under \texttt{src/} and integrate background global hotkey listener via \texttt{pynput}.
    \item \textbf{Assignee}: \texttt{Ysn-Ir}
\end{itemize}

\begin{terminalbox}
$ git checkout -b feature/architecture-and-hotkeys
Switched to a new branch 'feature/architecture-and-hotkeys'

$ git add .
$ git commit -m "feat(core): design project architecture and global hotkeys listener (fixes #1)"
[feature/architecture-and-hotkeys 4d4966e] feat(core): design project architecture and global hotkeys listener (fixes #1)
 10 files changed, 381 insertions(+), 11 deletions(-)
 create mode 100644 .gitignore
 create mode 100644 docs/GIT_WORKFLOW_LOG.md
 create mode 100644 src/capture/__init__.py
 create mode 100644 src/capture/base.py
 create mode 100644 src/config/__init__.py
 create mode 100644 src/config/settings.py
 create mode 100644 src/core/__init__.py
 create mode 100644 src/core/hotkeys.py

$ git push -u origin feature/architecture-and-hotkeys
branch 'feature/architecture-and-hotkeys' set up to track 'origin/feature/architecture-and-hotkeys'.
To https://github.com/hachemite/recordia-ilia.git
 * [new branch]      feature/architecture-and-hotkeys -> feature/architecture-and-hotkeys

$ git checkout main
$ git merge --no-ff feature/architecture-and-hotkeys -m "Merge branch 'feature/architecture-and-hotkeys' (Closes #1)"
$ git push origin main
Merge made by the 'ort' strategy.
To https://github.com/hachemite/recordia-ilia.git
   f0d65b2..d0735a3  main -> main
\end{terminalbox}

\begin{gitconcept}{Closing Issues via Merge Messages}
By including \texttt{Closes \#1} in the merge commit message, GitHub automatically transitions Issue \#1 to Closed when \texttt{main} is pushed.
\end{gitconcept}

\subsection{Step 2: Issues \#3 \& \#6 -- File Saving \& Settings Dialog}
\begin{itemize}[noitemsep]
    \item \textbf{Goal}: Build structured directory management for captures (\texttt{Screenshots/} and \texttt{Recordings/}) and an interactive GUI settings modal.
    \item \textbf{Assignee}: \texttt{Ysn-Ir} (covers unassigned duplicate \#6)
\end{itemize}

\begin{terminalbox}
$ git checkout -b feature/file-saving-and-settings
Switched to a new branch 'feature/file-saving-and-settings'

$ git add .
$ git commit -m "feat(storage): implement file saving manager and settings dialog (fixes #3, fixes #6)"
[feature/file-saving-and-settings e923205] feat(storage): implement file saving manager and settings dialog (fixes #3, fixes #6)
 6 files changed, 240 insertions(+), 19 deletions(-)
 create mode 100644 src/storage/__init__.py
 create mode 100644 src/storage/file_manager.py
 create mode 100644 src/ui/settings_window.py

$ git push -u origin feature/file-saving-and-settings
$ git checkout main
$ git merge --no-ff feature/file-saving-and-settings -m "Merge branch 'feature/file-saving-and-settings' (Closes #3, Closes #6)"
$ git push origin main
Merge made by the 'ort' strategy.
To https://github.com/hachemite/recordia-ilia.git
   d0735a3..e49e456  main -> main
\end{terminalbox}

\begin{notebox}{GitHub API Verification}
GitHub API inspection confirmed that both Issue \#3 and duplicate Issue \#6 were automatically closed upon push of merge commit \texttt{e49e456}.
\end{notebox}

\subsection{Step 3: Issue \#4 -- Screenshot Capture Engine}
\begin{itemize}[noitemsep]
    \item \textbf{Goal}: Build full-screen and region screenshot capture with resilient fallback for headless environments.
    \item \textbf{Assignee}: \texttt{Ysn-Ir}
\end{itemize}

\begin{terminalbox}
$ git checkout -b feature/screenshot-capture
$ git add .
$ git commit -m "feat(capture): implement screenshot capture engine with mss (fixes #4)"
[feature/screenshot-capture b3864cf] feat(capture): implement screenshot capture engine with mss (fixes #4)
 4 files changed, 156 insertions(+), 10 deletions(-)
 create mode 100644 src/capture/screenshot.py

$ git push -u origin feature/screenshot-capture
$ git checkout main
$ git merge --no-ff feature/screenshot-capture -m "Merge branch 'feature/screenshot-capture' (Closes #4)"
$ git push origin main
To https://github.com/hachemite/recordia-ilia.git
   e49e456..d2ab6c1  main -> main
\end{terminalbox}

\subsection{Step 4: Issue \#5 -- Screen Video Recording Engine}
\begin{itemize}[noitemsep]
    \item \textbf{Goal}: Implement multi-threaded desktop video capture using OpenCV \texttt{cv2.VideoWriter} (\texttt{mp4v} codec) and \texttt{mss}.
    \item \textbf{Assignee}: \texttt{Ysn-Ir}
\end{itemize}

\begin{terminalbox}
$ git checkout -b feature/screen-recording
$ git add .
$ git commit -m "feat(recorder): implement threaded screen recording engine with OpenCV (fixes #5)"
[feature/screen-recording ee15286] feat(recorder): implement threaded screen recording engine with OpenCV (fixes #5)
 4 files changed, 235 insertions(+), 14 deletions(-)
 create mode 100644 src/capture/recorder.py

$ git push -u origin feature/screen-recording
$ git checkout main
$ git merge --no-ff feature/screen-recording -m "Merge branch 'feature/screen-recording' (Closes #5)"
$ git push origin main
To https://github.com/hachemite/recordia-ilia.git
   d2ab6c1..1b0f505  main -> main
\end{terminalbox}

\subsection{Step 5: Issue \#7 -- Comprehensive Automated Test Suite}
\begin{itemize}[noitemsep]
    \item \textbf{Goal}: Construct unit and integration tests covering settings, file management, hotkeys, screenshot, and video recording.
    \item \textbf{Assignee}: \texttt{Ysn-Ir}
\end{itemize}

\begin{terminalbox}
$ git checkout -b feature/test-features
$ python -m unittest discover -s tests -v
test_directory_creation (test_file_manager.TestFileManager) ... ok
test_generate_recording_path (test_file_manager.TestFileManager) ... ok
test_generate_screenshot_path (test_file_manager.TestFileManager) ... ok
test_update_base_directory (test_file_manager.TestFileManager) ... ok
test_init_and_bindings (test_hotkeys.TestHotkeyManager) ... ok
test_register_hotkey (test_hotkeys.TestHotkeyManager) ... ok
test_stop_when_not_started (test_hotkeys.TestHotkeyManager) ... ok
test_start_and_stop_recording (test_recorder.TestScreenRecorder) ... ok
test_capture_fullscreen (test_screenshot.TestScreenshotEngine) ... ok
test_capture_region (test_screenshot.TestScreenshotEngine) ... ok
test_default_values (test_settings.TestSettings) ... ok
test_set_and_save_persistence (test_settings.TestSettings) ... ok
----------------------------------------------------------------------
Ran 12 tests in 0.976s - OK

$ git add .
$ git commit -m "test(core): implement comprehensive unit and integration test suite (fixes #7)"
$ git push -u origin feature/test-features
$ git checkout main
$ git merge --no-ff feature/test-features -m "Merge branch 'feature/test-features' (Closes #7)"
$ git push origin main
To https://github.com/hachemite/recordia-ilia.git
   1b0f505..798cf8a  main -> main
\end{terminalbox}

\subsection{Step 6: Issue \#8 -- User Documentation \& Setup Guide}
\begin{itemize}[noitemsep]
    \item \textbf{Goal}: Author comprehensive documentation, badges, architecture overview, dependency manifest (\texttt{requirements.txt}), and setup instructions.
\end{itemize}

\begin{terminalbox}
$ git checkout -b feature/documentation-and-report
$ git add .
$ git commit -m "docs: complete user documentation, setup guide, and Git workflow report (fixes #8)"
$ git push -u origin feature/documentation-and-report
$ git checkout main
$ git merge --no-ff feature/documentation-and-report -m "Merge branch 'feature/documentation-and-report' (Closes #8)"
$ git push origin main
To https://github.com/hachemite/recordia-ilia.git
   798cf8a..9f431ab  main -> main
\end{terminalbox}

\subsection{Step 7: Bonus 1 -- Automated Steps Recorder \& HTML Reporting}
\begin{itemize}[noitemsep]
    \item \textbf{Goal}: Transform Recordia into a Problem Steps Recorder that isolates each capture session into a project directory and compiles responsive HTML reports.
\end{itemize}

\begin{terminalbox}
$ git checkout -b feature/steps-recorder
$ git add .
$ git commit -m "feat(steps): implement automated steps recorder with window inspection and HTML reports"
$ git push -u origin feature/steps-recorder
$ git checkout main
$ git merge --no-ff feature/steps-recorder -m "Merge branch 'feature/steps-recorder' into main"
$ git push origin main
To https://github.com/hachemite/recordia-ilia.git
   9f431ab..bbece0a  main -> main
\end{terminalbox}

\subsection{Step 8: Bonus 2 -- Event-Driven Auto Step Triggers (Click \& Timer)}
\begin{itemize}[noitemsep]
    \item \textbf{Goal}: Provide users the choice between Manual capture, Auto on Mouse Click (captures every click with coordinates), and Auto Periodic Timer (captures every $N$ seconds).
\end{itemize}

\begin{terminalbox}
$ git checkout -b feature/auto-step-triggers
$ git add .
$ git commit -m "feat(triggers): implement automated step capture triggers for mouse click and periodic timer"
$ git push -u origin feature/auto-step-triggers
$ git checkout main
$ git merge --no-ff feature/auto-step-triggers -m "Merge branch 'feature/auto-step-triggers' into main"
$ git push origin main
To https://github.com/hachemite/recordia-ilia.git
   bbece0a..2ed5baa  main -> main
\end{terminalbox}

\subsection{Step 9: Bonus 3 -- Custom Hotkey Re-binding \& Step Description Editor}
\begin{itemize}[noitemsep]
    \item \textbf{Goal}: Allow users to configure their own hotkeys in Settings, review/edit step descriptions during or after sessions, and remove emoji stickers for a clean interface.
\end{itemize}

\begin{terminalbox}
$ git checkout -b feature/hotkey-config-and-step-editor
$ git add .
$ git commit -m "feat(ui): add custom hotkey configuration, step description editor, and clean sticker-free UI"
$ git push -u origin feature/hotkey-config-and-step-editor
$ git checkout main
$ git merge --no-ff feature/hotkey-config-and-step-editor -m "Merge branch 'feature/hotkey-config-and-step-editor' into main"
$ git push origin main
To https://github.com/hachemite/recordia-ilia.git
   2ed5baa..d94d961  main -> main
\end{terminalbox}

\vspace{0.8em}

% -------------------------------------------------------------
% QUALITY ASSURANCE & TEST SUITE
% -------------------------------------------------------------
\section{Quality Assurance \& Test Verification}

A robust automated test suite was maintained throughout the development lifecycle to ensure continuous regression protection across all modules:

\begin{terminalbox}
$ python -m unittest discover -s tests -v
test_init_and_state (test_auto_stepper.TestAutoStepper) ... ok
test_periodic_timer_trigger (test_auto_stepper.TestAutoStepper) ... ok
test_stop_when_idle (test_auto_stepper.TestAutoStepper) ... ok
test_directory_creation (test_file_manager.TestFileManager) ... ok
test_generate_recording_path (test_file_manager.TestFileManager) ... ok
test_generate_screenshot_path (test_file_manager.TestFileManager) ... ok
test_update_base_directory (test_file_manager.TestFileManager) ... ok
test_init_and_bindings (test_hotkeys.TestHotkeyManager) ... ok
test_register_hotkey (test_hotkeys.TestHotkeyManager) ... ok
test_stop_when_not_started (test_hotkeys.TestHotkeyManager) ... ok
test_start_and_stop_recording (test_recorder.TestScreenRecorder) ... ok
test_capture_fullscreen (test_screenshot.TestScreenshotEngine) ... ok
test_capture_region (test_screenshot.TestScreenshotEngine) ... ok
test_finish_session_generates_html_and_markdown_reports (test_session_manager.TestSessionManager) ... ok
test_record_steps_with_metadata (test_session_manager.TestSessionManager) ... ok
test_start_session_creates_directory_structure (test_session_manager.TestSessionManager) ... ok
test_update_step_description (test_session_manager.TestSessionManager) ... ok
test_default_values (test_settings.TestSettings) ... ok
test_set_and_save_persistence (test_settings.TestSettings) ... ok
test_get_active_window_info_structure (test_window_info.TestWindowInfo) ... ok
----------------------------------------------------------------------
Ran 20 tests in 2.579s
OK
\end{terminalbox}

\vspace{0.8em}

% -------------------------------------------------------------
% HOW TO RUN
% -------------------------------------------------------------
\section{How to Run \& Verify}

\begin{enumerate}
    \item \textbf{Clone the repository}:
    \begin{lstlisting}
$ git clone https://github.com/hachemite/recordia-ilia.git
$ cd recordia-ilia
    \end{lstlisting}
    \item \textbf{Install dependencies}:
    \begin{lstlisting}
$ pip install -r requirements.txt
    \end{lstlisting}
    \item \textbf{Execute the automated test suite}:
    \begin{lstlisting}
$ python -m unittest discover -s tests -v
    \end{lstlisting}
    \item \textbf{Launch the desktop application}:
    \begin{lstlisting}
$ python src/main.py
    \end{lstlisting}
\end{enumerate}

\vspace{0.8em}

% -------------------------------------------------------------
% CONCLUSION
% -------------------------------------------------------------
\section{Conclusion \& Collaborative Reflections}

The development of \textbf{RecordiaILIA} served as an exemplary practical demonstration of professional version control and agile collaboration:

\begin{enumerate}[leftmargin=*]
    \item \textbf{Conflict Avoidance}: By strictly separating features into dedicated branches and pulling upstream before merging, zero merge conflicts occurred throughout the project.
    \item \textbf{Transparent Audit Trail}: Non-fast-forward (\texttt{--no-ff}) merges ensure that even years after release, any team member or evaluator can visually trace the origin, testing, and integration of each individual feature in the Git commit graph.
    \item \textbf{Automated Governance}: Incorporating issue closure triggers (\texttt{Fixes \#X}) directly into version control messages harmonized the Git repository with GitHub's project management board, eliminating manual synchronization overhead.
    \item \textbf{Software Excellence}: Beyond Git mechanics, the resulting application is a complete, feature-rich desktop utility providing automated problem steps recording, customizable hotkeys, and multi-threaded screen recording.
\end{enumerate}

\end{document}
